%!TEX root = thesis.tex
\chapter{Synthesis and Cross-Paper Discussion}
\label{chap:ch5}

% TODO — introductory paragraph

\section{From Zone Definition to Zone Classification}
\label{sec:ch5:foundation}
% TODO

\section{What the 100\% Accuracy Result Means for the Zone Framework}
\label{sec:ch5:rf100}
% TODO

\section{Interpretability, Accuracy, and Deployability: A Three-Way Trade-Off}
\label{sec:ch5:tradeoff}
% TODO

\section{The Role of Physics in Machine Learning for PHM}
\label{sec:ch5:physics}
% TODO — this is the most important section for new citations.
% Cite \citep{gijon2024} here as a parallel example of physics-guided ML in wind energy:
% they combine a physics power-curve model with a neural network for residuals, and use
% SHAP for interpretability. Their conclusion matches this thesis: physics + data-driven
% is stronger than either alone.
% Also reinforce with \citep{hong2025} — physics constraints improve SCADA data quality,
% same principle as physics-informed zone boundaries improving classification quality.

\section{Practical Recommendations for Wind Farm Operators}
\label{sec:ch5:recommendations}
% TODO
